\section{Erd\H{o}s Problem 69: a prime-divisor series}

\subsection{1) Statement and current resolution}
Let $\omega(n)$ count the distinct prime divisors of $n$, with
$\omega(1)=0$. Is
\[
 S=\sum_{n\geq1}\frac{\omega(n)}{2^n}
\]
irrational? Tao and Ter\"av\"ainen proved this unconditionally in
Theorem 1.3 of their 2025 paper. The present attempt proves exact
identities, certified approximation bounds and a necessary rationality
condition. It does not reconstruct their correlation argument, so its
own proof is partial despite the established affirmative resolution.

\subsection{2) A convergent prime expansion}
The product of the distinct prime divisors of $n$ is at least
$2^{\omega(n)}$ and divides $n$. Thus $\omega(n)\leq\log_2 n$,
and $S$ converges absolutely. Since all terms are nonnegative, summing
first over prime divisors is justified and gives
\[
 S=\sum_p\sum_{k\geq1}2^{-kp}
  =\sum_p\frac1{2^p-1}.
\]
This is a sum over all multiples of a prime, not merely its powers:
confusing the two would replace $\omega$ by a different arithmetic
function. For any integer $P\geq2$, the prime truncation has the
rigorous error bound
\[
 0<S-\sum_{p\leq P}\frac1{2^p-1}
 \leq 2\sum_{m>P}2^{-m}=2^{1-P}.
\]
The strict positivity follows from the infinitude of primes.

\subsection{3) What rationality would force}
Suppose $S=a/b$, with integers $a,b>0$. For each $N\geq1$, set
\[
 T_N=\sum_{h\geq1}\frac{\omega(N+h)}{2^h}.
\]
Multiplying the original series by $2^N$ gives
\[
 T_N=2^N\frac ab-\sum_{n=1}^N\omega(n)2^{N-n}
       \in \frac1b\mathbb Z.
\]
In particular every $T_N$ belongs to one fixed discrete lattice,
although it need not be bounded. This is a useful exact obstruction
for an irrationality proof; a mere large value of $\omega(N+h)$ does
not contradict it.

Truncation can be controlled uniformly. Concavity of logarithm gives
\[
 \log_2(N+h)\leq\log_2 N+\frac{h}{N\log 2}.
\]
Using $\sum_{h>H}2^{-h}=2^{-H}$ and
$\sum_{h>H}h2^{-h}=(H+2)2^{-H}$, we obtain
\[
 0\leq T_N-\sum_{h=1}^H\frac{\omega(N+h)}{2^h}
 \leq 2^{-H}\left(\log_2N+\frac{H+2}{N\log2}\right).
\]
Thus rationality implies that every displayed finite weighted block
is within this explicit error of $b^{-1}\mathbb Z$. It would suffice
to find, for every fixed $b$, some $N,H$ whose weighted block has
distance from that lattice larger than the error bound.

\subsection{4) The remaining gap}
No argument here produces such blocks for every denominator $b$.
The coefficients $\omega(n)$ are not binary digits: they exceed $1$,
so binary carries prevent an inference from nonperiodicity of
$\omega(n)$ to irrationality of $S$. Likewise, arbitrarily accurate
rational approximations with uncontrolled denominators prove nothing
about irrationality. The known solution obtains the needed control
using quantitative correlations of multiplicative functions; that
external proof is cited for the problem's status, not silently counted
as a proof supplied by this attempt.

\textbf{Outcome of this attempt: UNRESOLVED; exact reductions and
error estimates proved.} No new mathematical advance is claimed.
The batch verification checks the divisor/prime rearrangement on
finite sums using exact rational arithmetic.

\subsection{5) Sources}
\begin{itemize}
\item T. Tao and J. Ter\"av\"ainen, \emph{Quantitative correlations
and some problems on prime factors of consecutive integers}, Theorem
1.3 and Section 5,
\texttt{https://arxiv.org/abs/2512.01739}.
\item T. F. Bloom, Problem 69,
\texttt{https://www.erdosproblems.com/69}, checked 23 September 2026.
\end{itemize}
The percentage is a conservative estimate of the proof supplied here,
not of the confidence in the published resolution.
\subsection{6) COMPLETION ESTIMATE}
\textbf{COMPLETION: 10\%}
